% Records document: detail moved verbatim from the posted note in the 2026-11 compression (DECISION-20261119c-joseph-note-compression-outline.md
% item 3; move-manifest-je JE-10, JE-12). It holds two kinds of text: the history of stopped tests and superseded comparisons
% (class H, the director's App. E ruling), and current detail of later reporting-only studies that the compressed note states
% in short form. Neither is a struck value; app_history.tex keeps those.
\section{Detail moved from the note in the compression}
\label{app:compression-records}

% Moved verbatim from the note's development-history appendix in the compression (DECISION-20261119c-joseph-note-compression-outline.md
% item 3; the director's App. E ruling of 2026-11-19: class H, stopped tests with no result and superseded comparisons; move-manifest-je JE-10).
\subsection{Development-history detail moved from the note}

\subsubsection{The earlier two-dimensional budget: comparison}
\paragraph{Comparison with the earlier 2D budget.} The saved products and code of the earlier 2D budget~\cite{TwoDFollowup2026} show the following defects in its construction; the quoted budget corrects each. First, its systematic sweep has a simulated truth input that differs from that of its seed-42 reference, while its saved reconstructed inputs agree. A common displacement cancels to first order in the MAT mean-centered covariance: the mismatch does not invalidate the entire covariance and supplies no correction to subtract, and because the earlier input and code revision are not available, it cannot be attributed more specifically. The sweep used by the quoted universe block does not have this mismatch.
% (one paragraph since the history restructure's style pass, 2026-11-15: history-cuts/sec_systematics.md S6)
Second, the bands Rvn1pi and Rvp1pi encode the same shift under the configured nonresonant-pion tune, and the earlier budget counted it twice; the quoted universe block drops Rvp1pi and counts it once (\S\ref{sec:uq-universes}). The band is small, \SI{0.088}{\percent}, but the double count is a defect in the construction of the covariance; no conclusion about the published analysis's implementation follows.
%
Third, the earlier budget did not propagate background-template statistics, whose error exceeds one tenth of the data's Poisson error in 60 of 205 saved reconstructed bins (a reconstructed-space diagnostic, not a measured unfolded uncertainty); the quoted statistical block draws the background simulation as its own bootstrap stream (\S\ref{sec:uq-boot}). The values printed for the earlier budget were not changed by these findings.

\subsubsection{The stopped coverage tests of the sixteen-cell covariance: design and stops}
A first coverage test at these two tunes was specified before any pseudo-experiment was generated. It tested the covariance without the run-period block, with pseudo-experiments that also draw the flux, interaction and detector uncertainties (the lateral and normalization ones as surrogates, and $C_{\mathrm{model}}$ and $C_{\mathrm{resp}}$ not at all). It was stopped before any coverage pseudo-experiment ran. App.~\ref{sec:cov1b-gate} gives the result at one of its gates and the investigation of that gate. Replacing the pseudo-experiments' values by the covariance's for any \fvval{cov1b-kmax-word} or fewer of the \fvval{cov1b-nbands} drawn bands does not bring every cell inside that gate's window.
% F16-2 (2026-11-16): the tests numbered as the paper and Supplement number them (paper_body.tex "The first five-observable
% test"; supplemental_body.tex S8: "A first coverage test ... A staged redesign then stopped at its first stage"): this gate-
% stopped test is the first, the staged run below is its redesign, and LL5D is the second; App. I's two rows say the same.

A staged redesign of that first test, at the two tunes and with the covariance rebuilt by the same procedure at each tune, also ended without a coverage result. Each of its pseudo-experiments is drawn from a simulated world built from randomly drawn systematic universes; in each cell, that world's denominator is the first drawn universe's denominator plus each later universe's departure from the central-value denominator. A rule fixed before the first stage ran stops that stage if any draw gives a non-positive denominator in a cell that is positive at the central value. \lrval{stop-flagged} of the \lrval{stop-draws} draws did so, with at most \lrval{stop-maxcells} such cells in a draw, and the first stage stopped. The record does not establish why those cells are non-positive.
% absence-ok: "The record does not establish why those cells are non-positive" is the source paragraph's own statement (sec_fived16.tex:283 @ 2fe070fec, LR-12), moved verbatim

\subsubsection{Lost pseudo-experiments of the generator tests: evaluation of the recovered draws}
The recovered draws were evaluated in two ways: with the variant shifts of the primary evaluation held fixed (the \emph{fixed-shift} evaluation), and with the shifts recomputed on all draws (the \emph{recomputed-shift} evaluation). The recovery's checks are narrower in scope than the main recomputation. Only the fixed-shift evaluation was cross-checked independently. The recomputed-shift evaluation was reproduced only by a standalone rerun of the same code, a consistency check rather than an independent recomputation. The earlier stops rest on the calibration run's own re-application of the sequential rule and have no independent check.

\subsubsection{Investigation of the gate of the sixteen-cell coverage test}
\label{sec:cov1b-gate}
% History restructure (2026-11-14): moved verbatim from sec_fived16.tex (sec:fived16-fd) at 1fcb3713d:281-285, with its
% source comments 1fcb3713d:292-296; history-cuts/sec_fived16.md row G1. The test and its stop are in sec:fived16-fd.
The coverage test at the two tunes for the covariance without the run-period block (\S\ref{sec:fived16-fd}) was stopped before any coverage pseudo-experiment ran. At one of its gates, the ratio of each cell's uncertainty reproduced through the pseudo-experiments to that of the drawn covariance lay outside \fvval{cov1b-win-lo}--\fvval{cov1b-win-hi} in \fvval{cov1b-ncells-out} of the 16 cells, ranging from \fvval{cov1b-rho-min} to \fvval{cov1b-rho-max}. This subsection gives the investigation of that gate.

For the interaction uncertainties, a later investigation at the nominal truth found that the gate compared unlike objects. Its pseudo-experiments reproduce in every cell the spread of the response alone, with the simulated truth held fixed, whereas the full interaction universes from which the covariance is built give a different spread at the same truth: the ratio of the response-only spread to the full-universe spread has median \lrval{i2-univ-med} (range \lrval{i2-univ-lo}--\lrval{i2-univ-hi}). As a sensitivity, the standard deviation of the response alone, carried into the data's frame, is a median \lrval{i2-resp-med} times that of the reported interaction block (range \lrval{i2-resp-lo}--\lrval{i2-resp-hi} over the cells); it depends on the fine-cell grid and is not a bound. Separately, for the interaction uncertainties, in all but one cell the covariance's construction gives a different spread at the nominal truth than at the data, where the covariance is made. Pseudo-data that take each full interaction universe as the truth, unfolded with the nominal simulation, give a spread whose ratio to that of the covariance's construction at the nominal truth lies within the gate's window, \fvval{cov1b-win-lo}--\fvval{cov1b-win-hi}, in all 16 cells: under this specified swap the universes' second-order departures at $\pm1\sigma$ are below the window, as expected, since closure fixes the ratio at one to first order.

This concerns the test; it neither validates nor undermines the reported interaction block, and neither finding shows which spread the reported intervals should carry. In cells R00 and R05 the interaction uncertainties carry only 0.57 and 0.36 of the squared ratio's departure from one, so this account does not explain the gate's departure there, and this account says nothing about the parts carried by the flux and detector uncertainties.

A later pre-check of the flux and detector groups at the nominal truth found that, as for the interaction group, the pseudo-experiments reproduce, within the gate's window in every resolved cell, the block that the same weight factors give on the simulation side. Within the detector group, the two sides give blocks outside that window for the GEANT neutron band in cells \lrval{p1-neutron-cells} and for the GEANT proton band in cells \lrval{p1-proton-cells}, with no cause established. A diagnosis made afterward found the two sides' inputs equal and the departures stable across estimator seeds; its test of how the departures scale with the band's size reads mixed in both bands, so no mechanism is established.
% U12 (currency check of 2fe070fec): OUTCOME-20261109-band-diagnosis.md (commit de5fc4676): inputs EQUAL (P2), SEED-STABLE in all 10 departing cells (P4), scale MIXED in both bands (P4), "no mechanism is licensed"; readout fd/receipts-1b/band-diag/ (553c3dff).
% absence-ok: "no mechanism is established" restates the record's own verdict, "no mechanism is licensed" (OUTCOME-20261109-band-diagnosis.md title and readings table)
% Note round 1, N3 (director's go, 2026-11-08): the response-only sensitivity, DECISION-20261101-joseph-1b-keep-block-redesign-test.md
% :38-40 @ b551b05f ("goes to the note as a sensitivity"; grid-dependent, not a bound): sqrt(B/M) per cell from
% fd/receipts-1b/inv/i2/i2.json (a9d3dc9c), checked against FD-1B-INV-I2-REPORT.md:299 and printed at the decision's precision.
% Revision C' 2026-11-06 (two referee reads of C, Sonnet and Grok, both REVISE; the director's five changes): the R00/R05 caveat, M at the data cited, kappa as stage 0 states it, the window named, "fine cell" dropped, and the test-not-covariance clause. Records (docs/orchestration/state/measurement-paper/5d-community-20261024/, all @ ddc35ae15 unless stated): FD-1B-INV-I2-REPORT.md:13-15 @ 5da29ca4 (O1, unlike objects = B against A; F_src = 0.986-1.015 in every cell) and :131 ("M is S1's, made at the data. The window is [0.9, 1.1]"), sec. 4.1 :187 (F_truth = sqrt(A/M), a separate truth-dependence factor, outside the window in 15 of 16) and :191 (S1's mechanism at the nominal truth gives about 0.68x S1's sigma at the data), sec. 4.3 :214 (F_prior outside in all 16); FD-1B-INV-DECLARATION.md:107 (S1's blocks are made at the data and gate D's at the nominal truth); FD-1B-MODED-DECLARATION.md:95 ("M is made at (real data, m0)") and :416 ("The reported intervals are made with M at the data's point"); FD-1B-INV-I2-DIRECTOR-REVIEW.md:39-40 ("O1 is a statement about the test. It neither validates nor undermines the reported block") and :41-42 (R00 and R05: the interaction group carries only 0.57 and 0.36 of rho^2-1, not explained); FD-1B-PRECHECK-ATTRIBUTION.md:82,84 (interaction shares 0.565 and 0.359; flux 0.410 and 0.576), :56-57 (Dnom: full-universe pseudo-data at the nominal truth, unfolded with the nominal simulation), NOT CITABLE FOR :19 (no harness statement about flux or detector); FD-1B-STAGED-STAGE0-REPORT.md:12-16 @ a97dfa47 (kappa = sqrt(Dnom/A_unc) in [0.9, 1.1] in all 16; the full universe's second-order departures under the swap are below the window; closure fixes kappa = 1 at first order, the expected outcome); fd/receipts-1b/stage2/run/score/gate-d.json:53-56 (window [0.9, 1.1], the same as kappa's) and :114-115 @ 1a4815a8. Revision C''/E'' 2026-11-06 (confirmation referee read, minor REVISE): flux/detector sentence reworded so it makes no claim about the gate, whose ratio includes flux (FD-1B-PRECHECK-ATTRIBUTION.md:82,84: flux 0.410 in R00, 0.576 in R05); the Separately sentence scoped to the interaction group (I2 :187 counts the interaction group); "reported interaction block" as FD-1B-INV-I2-DIRECTOR-REVIEW.md:39-40 words it; departures below the window at +-1 sigma (FD-1B-STAGED-STAGE0-REPORT.md:14-15 @ a97dfa47); kappa's denominator A_unc = A / g^2 (FD-1B-PRECHECK-ATTRIBUTION.md:57; FD-1B-MODED-DECLARATION.md:70), with A S1's whole universe at the nominal truth (FD-1B-PRECHECK-ATTRIBUTION.md:54-55, citing FD-1B-INV-DECLARATION.md:225-229), i.e. the covariance's construction at the nominal truth without the carry g.
% Director-added correction 2026-11-06, for Joseph's review before merge (not among the two he approved): "why the ratios depart from one has not been established" was stale for the interaction group at the nominal truth. INV I2 (#650): docs/orchestration/state/measurement-paper/5d-community-20261024/FD-1B-INV-I2-REPORT.md:13-15 @ 5da29ca4 (O1, "gate D compared unlike objects"; F_src = sqrt(E/B) = 0.986-1.015 in every cell; sqrt(B/A) outside the window in all 16, sec. 4.3; F_truth = sqrt(A/M) outside in 15 of 16, sec. 4.1), quoted with FD-1B-INV-I2-DIRECTOR-REVIEW.md sec. 3's narrowed labels (no "prior part"; O1 neither validates nor undermines the reported block). Staged stage 0 (#667): docs/orchestration/state/measurement-paper/5d-community-20261024/FD-1B-STAGED-STAGE0-REPORT.md:12-13 @ a97dfa47 (kappa = sqrt(Dnom/A_unc) in [0.9, 1.1] in 16 of 16; Dnom = pseudo-data at means with each of the 67 interaction universes as nature; read at the nominal truth, not M). Gate D: docs/orchestration/state/measurement-paper/5d-community-20261024/fd/receipts-1b/stage2/run/score/gate-d.json:114-115 @ 1a4815a8 (response only, truth kept in every fine cell, against S1's blocks). CANNOT SUPPORT (both reports): coverage in either direction; which block is right; why A and B differ; the flux and detector groups; gate D in R00 and R05, where the interaction group carries 0.565 and 0.359 of rho^2 - 1.

% Move pair H3-020 (director-helper-3; move-manifest-h3-20261117.tsv sends): sec_results.tex:148 @ 79a934a41, sentences 2-4,
% verbatim with the source comments :153-154; sec:results-tension keeps sentence 1 and points here.
\subsubsection{The exact-split cross-check's eigenmode decomposition (from \S\ref{sec:results})}
\label{app:devhistory-results}

This is the exact-split cross-check's part of the eigenmode decomposition of \S\ref{sec:results-tension}. For the exact-split cross-check, $\chi^{2}$ is larger and distributed differently: its components of magnitude $5$--$6$ lie in moderate-$\lambda$ modes on the low-$\pT$ ridge of the cross-section peak. Its 10 smallest-$\lambda$ modes carry \SI{3}{\percent} of the $\chi^{2}$, and the 73 best-measured directions give $\chi^{2}=291$ ($\chi^{2}/205=1.42$, again divided by the number of reported bins; App.~\ref{sec:tension}). Going from 5 to 10 iterations changes the LightGBM (seed~1) value by a negligible $+0.01$.
% Revision D' 2026-11-06 (referee reads of D, both ACCEPT; Sonnet and the interpreter: print the raw chi2 and the directions kept): docs/orchestration/receipts/2d-lgbm42-20261011/tension_report_exact_control.txt:56 prints "1e-04  1.419  73". Arithmetic: 1.419 x 205 = 290.9; 205 x [1.4185, 1.4195) is in [290.8, 291.0], 291 at three figures.
% Director-added wording correction 2026-11-06, for Joseph's review before merge (not among the two he approved): the truncated chi2 was printed as chi2/ndf, but its divisor is the 205 reported bins, not the directions kept. 2d-unfolding/diagnose_tension.py @ c3dbe30f4 (main): :110 n = int(mask.sum()), the 205 reported bins; :183 and :193 divide every truncated chi2 by n (the same divisor at 56e30b9c9 :108/:181/:191 and at 735c909da). Receipt docs/orchestration/receipts/2d-lgbm42-20261011/tension_report_exact_control.txt:56 ("1e-04  1.419  73"). No number changed.

% Moved verbatim (compression K9, class H: a superseded configuration's detail; move-manifest-je JE-10.028..).

% Moved verbatim from the note's generator tests (sec_joint.tex) in the compression (G2b; DECISION-20261119c item 3 and the director's rule (5),
% 2026-11-19: the later reporting-only studies' detail; the note keeps each study's conclusion and qualifiers; move-manifest-je JE-12).
\subsection{Generator-test detail moved from the note}

The multiplier $\kappa$ scales the last measured refinement change $\delta_{\mathrm{fine}}$: $\kappa=1$ means one such change, not one standard deviation or a confidence level. A scan of $\kappa$ along the direction $\delta_{\mathrm{fine}}$ of the claim variants~\cite{JointTestKappa2026} (for reporting only) found, in the calibration ensembles used for the decisions, all ten decisions to remain determinate rejections at every sampled multiplier up to 5.25. Determinate rejection is first lost between the sampled multipliers 5.2547 and 5.2563, and between 5.1688 and 5.1703 with the recovered draws included. It is lost by the GENIE~2.12.10 + MEC shape test, whose interval then contains its Holm threshold, and, through the Holm step-down, by the GENIE~2.12.10 shape test. Neither is then a determinate failure to reject. The scan probes one prescribed direction. It establishes neither convergence nor a bound on the actual below-grid residual (condition~(i), \S\ref{sec:joint5d-conditions}).
% Search: below the fine grid; sub-fine-grid; fine-grid convergence; delta_fine; kappa breakdown @ ddf16f10

Because every test rejects, power serves here only as context: the test specification uses power only to interpret a failure to reject. At the Tune~v1 null, two of the three specified alternatives are detected in shape with power 1.0 at the 0.005 level, and the third, a MEC-like $(\Eavail,W)$ shape, with \jtPowPthree. At the GENIE~2.12.10 null the shape powers at 0.005 are \jtPowPonegShape, 0.02 and 0.00. The binomial 95\% intervals are $[0.98,1]$ for the first two Tune~v1 alternatives, $[0.17,0.30]$ for the third, and $[0.67,0.80]$, $[0.01,0.05]$ and $[0,0.02]$ at the GENIE~2.12.10 null. These values use the power sets as evaluated for the decisions. Each of those sets had lost some of its 200 alternative pseudo-experiments to time limits (172--199 retained). Evaluated against the null ensembles used for the decisions, the complete power sets, recovered as in \S\ref{sec:joint5d-lost}, agree with them within 0.01. When the GENIE~2.12.10 null is also recomputed on all draws, the shape power at 0.005 against the first alternative at that null is 0.72.
% X7 (sol minor), moved from the article: docs/orchestration/state/s5p/stage7/joint/joint-evaluate.json
% power.<P>.shape["0.005"].claim_rule_determined.interval (n 193, 195, 195; 199, 193, 172).

The published comparison in muon kinematics alone, whose observables do not use the reconstructed recoil energy, also disfavors the same generator families~\cite{MINERvA:2021owq}. It provides context for a disagreement involving muon observables, but the hypotheses tested here use a different estimator and calibration.

In an exploratory sensitivity study of the response~\cite{JointTestRecoil2026}, the data are re-unfolded with the simulation's reconstructed recoil energy scaled by $\pm\SI{4}{\percent}$, and the ten tests are re-evaluated against the null distributions used for the decisions. This scale is the level of data--simulation agreement reported for test-beam protons, pions and electrons of 0.35--\SI{2.0}{GeV/c}~\cite{MINERvA:2015testbeam}. Neutrons and higher-momentum hadrons, which matter at high $W$, lie outside that sample. The $\pm4\%$ scale is not a validated uncertainty prescription. A re-unfold with unscaled recoil reproduces the decisions of Table~\ref{tab:joint5d}, with every claim $p$ value inside the range spanned by the data's rounding variations. At both $-4\%$ and $+4\%$ every one of the ten rejections holds. The observed statistics move by up to about \SI{11}{\percent} from their recorded values. Nine tests keep $k=0$. NuWro's shape test, the closest to its null, changes from $k=1$ to $k=3$ of 1751 at $+4\%$ ($p=0.0023$, against a Holm threshold of 0.05). The evaluation was reproduced independently. Table~\ref{tab:joint5d-recoil} gives the observed statistics.
% Search: validated uncertainty prescription; calorimetric response band; recoil-energy-scale band @ ddf16f10

% Deposit pass: W2B-REPORT-20261006.md sec. 5, "Observed statistics" (descriptive) and the k/B table; values copied as recorded. Null medians are the frozen unshifted (c = 0) null medians.
\begin{table}[htbp]
  \caption[Descriptive recoil-scale sensitivity of the observed statistics]{All ten tests remain rejected when the simulation's reconstructed recoil energy is scaled by $\pm4\%$ against unchanged null distributions. The observed joint statistics are shown beside the recorded values and the median of each unshifted null distribution, which is not regenerated. Every test keeps $k=0$ at both signs except NuWro's shape test ($k=1$ of 1751 at $-4\%$, 3 at $+4\%$). The values are descriptive.}
  \label{tab:joint5d-recoil}
  \centering\small
  \begin{tabular}{@{}lrrrrrrrr@{}}
    \toprule
    & \multicolumn{4}{c}{$T_{\mathrm{tot}}$} & \multicolumn{4}{c}{$T_{\mathrm{sh}}$} \\
    \cmidrule(lr){2-5}\cmidrule(lr){6-9}
    $G$ & recorded & $-4\%$ & $+4\%$ & null median & recorded & $-4\%$ & $+4\%$ & null median \\
    \midrule
    MINERvA Tune~v1     & 2852 & 2856 & 2540 & 49   & 631  & 647  & 600  & 49 \\
    GENIE~2.12.10       & 1813 & 1796 & 1721 & 528  & 1086 & 1104 & 1040 & 582 \\
    GENIE~2.12.10 + MEC & 1647 & 1652 & 1595 & 406  & 869  & 893  & 839  & 450 \\
    NuWro~21.09         & 7427 & 7380 & 7361 & 6324 & 9660 & 9690 & 9540 & 7110 \\
    GiBUU~2019          & 2867 & 3129 & 2665 & 1791 & 4013 & 4154 & 3901 & 1975 \\
    \bottomrule
  \end{tabular}
\end{table}

The study measures the sensitivity of the observed statistics and decisions to one coherent scale of reconstructed available energy and recoil energy, applied after the energy calibration at a single magnitude. It leaves untested (a) the response of the null distributions, which were not regenerated with the variation although they are simulated through the same detector model, so the size of this omission is unmeasured for these strongly non-central nulls; (b) particle-species, energy-dependent or resolution effects, and a response difference applied before calibration; and (c) larger variations. The rejections are therefore conditional on the simulation's hadronic response, and the response model is not shown to be complete.
% Register round n1 (2026-10-25): this subsection is the one home of the recoil-response omission; the condition
% (ii) gloss in sec:joint5d-conditions points here. The fixed null distributions print in the paragraph above and
% in (a); the need to vary them in a response-uncertainty construction, in the first paragraph.
% Search: regenerated with the; response-uncertainty; calorimetric response band; net effect on each p value @ ddf16f10

NuWro~21.09 and GiBUU~2019 were generated on carbon, while the data and the other predictions are per nucleon of CH. This study corrects the two predictions to CH and re-evaluates their tests against the null distributions used for the decisions. These target-corrected evaluations are descriptive. The decisions of Table~\ref{tab:joint5d} are unchanged. The primary generator comparisons (\S\ref{sec:fived16-gen}) do not use this correction and show NuWro and GiBUU on carbon only. This section records the correction.
% Descriptive, decisions unchanged: DECISION-20261020-joseph-nuwro-ch-claim.md:37 ("the replays are descriptive and
% do not re-grade A7"), reworded on the director's ruling of 2026-10-25 (register round n1).

\paragraph{Construction.} Each prediction is rebuilt per fine cell as $x_{\mathrm{CH}}=\tfrac{12}{13}\,x_{\mathrm{C}}+\tfrac{1}{13}\,x_{\mathrm{H}}$. Here $x_{\mathrm{H}}$ is the generator's own free-proton product, run through the unchanged conversion, flux-repair and binning chain. Each prediction's Monte Carlo variance is combined with the squares of the same weights. The hydrogen fraction is $f_{\mathrm{H}}=5.08\%$ for NuWro over the grid and 4.77\% for GiBUU over its 72 cells. Over the reported cells, the ratio $r_{c}=\mu_{\mathrm{CH}}/\mu_{\mathrm{C}}$ has minimum, median, and maximum 0.923, 0.963, and 1.039 for NuWro (109 cells) and 0.923, 0.963, and 1.036 for GiBUU (72 cells). Per cell, $f_{\mathrm{H}}$ reaches 0.112 and 0.109.  On these cells both corrections move the prediction toward the data. Its effect on the sixteen-cell comparison is stated in \S\ref{sec:fived16-gen-target}.
% B:29-34 (combination, sumw2, the frozen chain, the supplement); C:46-51 (GiBUU, frozen r(E), 72 cells).
% Accuracy fix (unscored; director 2026-10-26): "with the same weights" misstated the method. The code fills each
% event with weight c = w r / N x (12/13 or 1/13) and sums c^2 for the variance (docs/orchestration/state/
% measurement-paper/phase1/mp2/s2-receipts/d4/harmonize_9dbcde93.py:71, :104-108, :312, :418, :426 @ 7ad668f4), so
% sumw2_CH = (12/13)^2 sumw2_C12 + (1/13)^2 sumw2_H (mp3b/REPORT.md:30).
% Draft 6 (REVIEW-2 finding 12): B:97 (grid integral; 0.0506 over the 109 cells, B:105), C:51 (72-cell domain).
% B:97, :103-105; C:51, :56-58.
% B:138-139; C:80-81 (the cross term is negative: the data lie on the CH side). REVIEW-2 finding 11: pointer moved here.

\paragraph{Re-evaluated tests.} With the metric used for the decisions, NuWro's total statistic falls by 3.83 null standard deviations and its shape statistic by 0.37.  At its original Holm position (threshold 0.005) the total test would then lie 1.11 null standard deviations below its boundary, and so be undetermined.  It remains rejected only because its larger $p$ value moves it to the last Holm position (threshold 0.05). In the $\kappa=\jtMoneKappaRobust$ family, which is for reporting only, it is not rejected ($p=0.075$). NuWro's shape test stays rejected, 0.74 null standard deviations above its decision boundary.  GiBUU's statistics fall by 4.67 (total) and 1.33 (shape) null standard deviations. Both tests stay rejected at their original Holm positions, including at $\kappa=\jtMoneKappaRobust$. Their margins shrink from 7.86 to 3.18 and from 3.58 to 2.24 null standard deviations.  With each CH prediction's own variance in the metric, no decision changes for either generator.
% B:123-137 (NuWro: T_total 7427.1 -> 6759.8, -3.83 null SD; T_shape 9660.0 -> 9333.1, -0.37; frozen position:
% threshold 0.005, T_crit 6953.94, -1.11 SD, undetermined; re-sorted to the last position, threshold 0.05, CP
% [0.0183, 0.0336], rejected; shape to threshold 0.025, 0.74 SD above its frozen T_crit; kappa = 3: total p 0.0748,
% k = 130, not rejected; shape rejected). C:21, :72-79 (GiBUU: T_total -4.67, T_shape -1.33 null SD; margins 7.86 ->
% 3.18 and 3.58 -> 2.24; frozen positions, no re-sort; kappa = 3 also rejected). M:65-78 (frozen margins).
% B:129 (W_CH: both rejected, no change); C:76.

\paragraph{Generation checks.} For GiBUU, the sampling-closure check failed on the hydrogen sample. A bias check of that check's premise was added post hoc. Its criterion, a change of at most $0.1\,\sigma_{c}$ in every cell, was specified before the check was computed. The check met the criterion: it moved the CH prediction by at most $5.1\times10^{-4}\,\sigma_{c}$. Here $\sigma_{c}$ is the standard deviation of the cell's result over NuWro's null distribution used for the decisions ($B=1751$). The GiBUU re-evaluation is reported under the label ``B3 (sampling closure) failed on hydrogen; a post-hoc bias check of its premise passed at $\le0.1\,\sigma_{c}$''.

It has a further limitation: GiBUU's deep-inelastic events on hydrogen, 46\% of its hydrogen sample, do not conserve energy.  They lose a median of \SI{45}{MeV}, with a median transverse-momentum imbalance of \SI{31}{MeV}, while its other hydrogen events balance exactly. The cause has not been identified. The available energy of those events may be low, and the non-conservation may also affect GiBUU's carbon prediction, which cannot be checked in this way.

In NuWro's hydrogen generation, four chunks aborted and were rerun on alternate seeds. The fraction of events this excludes is estimated as $1.8\times10^{-5}$ in the main sample and $3.8\times10^{-5}$ in the 50--\SI{100}{GeV} supplement, or $1.9\times10^{-5}$ and $6.2\times10^{-5}$ including the first attempt. These estimates rest on four to six events.  The supplement's 95\% upper bound of $1.8\times10^{-4}$ exceeds the $10^{-4}$ level set for it, so the supplement is flagged rather than assumed negligible.
% C:7-8 (the required label, verbatim), :23-44 (declaration 00efbb25..., recorded before computing; max |Delta_c/sigma_c|
% 5.1e-4 against the 0.1 acceptance). B:48-56 (NuWro hydrogen chunks: 3 + 1 aborts rerun on alternate seeds; excluded
% fraction 1.8e-5 main, 3.8e-5 supplement; the supplement's 95% bound 1.8e-4 above the 1e-4 bar, FLAGGED).
% sigma_c: mp3c/REPORT.md:31 takes sigma_c from MP3-B; mp3b/REPORT.md:37, reading (a), "the SD of the J-cell result
% over NuWro's frozen B = 1751 null"; max |Delta_c/sigma_c| 5.1e-4 in reading (a), mp3c/REPORT.md:38.
% Draft 6: docs/orchestration/state/measurement-paper/BINDING-PHASE1-ADDENDUM-9.md:28-31 (sec. 4, "Recorded for the record";
% in tree since the rebase onto 51da2a62): production_ID 34, 46% of the GiBUU-H sample; median energy loss
% 45 MeV, median p_T imbalance 31 MeV; non-DIS hydrogen events balance exactly; cause not identified.
% ADDENDUM-9 :29-31 ("Carbon cannot be checked this way"; "possibly on the frozen GiBUU product").
% Search: do not conserve energy; energy non-conservation; production_ID 34 (cause not identified) @ ddf16f10

% Moved verbatim from the note's App. B (app_negweight.tex) in the compression (B1; move-manifest-je JE-13).
\subsection{Background-subtraction detail moved from the note}

\paragraph{Binned limit: the per-bin derivation (from App. B, sec:negweight-binned).} A step-1 classifier restricted to functions piecewise-constant on reconstructed analysis bins $\{b\}$ gives the population-level binned limit of the OmniFold/Richardson--Lucy correspondence in Refs.~\cite{Andreassen:2019cjw,Andreassen:2021zzk}. Its output is a single $p$ per bin. With common support and unnormalized class weights, step~1 weighted cross-entropy separates into per-bin objectives
\begin{equation}
  \label{eq:negweight-Lb}
  L_b(p) = -W_1\,\ln p - W_0\,\ln(1-p),
\end{equation}
where $W_0=S(b)>0$ is the integral of simulated class-0 density $S$ in bin $b$, and $W_1$ is the signed class-1 weight. Injecting background simulation at negative weight while keeping data at $+1$ gives $W_1=D(b)-B(b)$. Here $D(b)$ and $B(b)$ are the bin masses of $D$ and $B$, equivalently $N_{\mathrm{data}}(b)$ and $N_{\mathrm{bkg}}(b)$ in \S\ref{sec:negweight-estimator}.

For $D(b)>B(b)$, objective~\eqref{eq:negweight-Lb} is strictly convex on $(0,1)$ with unique interior minimizer $p=W_1/(W_0+W_1)$, giving
\begin{equation}
  % \label{eq:negweight-rb} is defined in the note (App. B keeps the equation); this copy drops the line.
  r(b) = \frac{p}{1-p} = \frac{W_1}{W_0} = \frac{D(b)-B(b)}{S(b)}.
\end{equation}

\paragraph{Binned limit: the divergence for D(b)<B(b) (from App. B).} The reduction fails where $D(b)<B(b)$. With $W_1<0$, $L_b(p)=|W_1|\,\ln p-W_0\,\ln(1-p)$ tends to $-\infty$ as $p\to0^+$ and is unbounded below, with no minimizer.


